%% kapa: exam=1 round=32 year=2021 period=1 form=A subject=civil_law subject_ko=민법 seq=17 number=17
%% source: Q-Net 감정평가사 기출문제 (cst00309) artlSeq=5208397 file=extracted/1008/5208397/A형/2021년도 제32회 감정평가사 1차 1교시 A형.pdf
%% conversion: kordoc 4.24.51 -> tools/gen1.py (auto); review=proofread-sonnet; flags=none
\begin{question}{17}
소급효가 원칙적으로 인정되지 않는 것은? (다툼이 있으면 판례에 따름)
\begin{choices}
  \item 무권대리인이 체결한 계약에 대한 추인의 효과
  \item 기한부 법률행위에서의 기한도래의 효과
  \item 토지거래 허가구역 내의 토지거래계약에 대한 허가의 효과
  \item 소멸시효 완성의 효과
  \item 법률행위 취소의 효과
\end{choices}
\end{question}
\answer{2}
